% META: {"id": "TCOMPL-MF-EX-012", "chapitre": "TCOMPL-MODELES-FONCTION", "type_objet": "exercice", "capacites": ["TCOMPL-MF-C5"], "capacites_codes": ["C5"], "methodes": ["M2"], "parcours": 2, "competences": ["calculer"], "duree_min": 15, "mode_creation": "ex_nihilo", "sources_inspiration": [], "fichier_tex": "chapitres/TCOMPL-MODELES-FONCTION/exercices/TCOMPL-MF-EX-012.tex", "corrige_tex": "chapitres/TCOMPL-MODELES-FONCTION/corriges/TCOMPL-MF-CO-012.tex", "status": "approved"}
% BEGIN-VERIFY
% from sympy import *
% x = symbols('x', positive=True)
% f = x*ln(x)
% fp = diff(f, x)
% assert simplify(fp) == ln(x) + 1
% assert solve(fp, x) == [exp(-1)]
% assert f.subs(x, exp(-1)) == -exp(-1)
% assert limit(f, x, oo) == oo
% assert limit(f, x, 0, '+') == 0
% END-VERIFY
\begin{exercice}{TCOMPL-MF-EX-012}{1}{15}
Etudier les variations de $f(x) = x\ln x$ sur $\left]0, +\infty\right[$, limites comprises.

\end{exercice}
